\chapter{为什么 Agent 需要“培育”而不仅是训练}
\label{chap:agent-cultivation-vs-training}

\begin{chapteroverviewbox}
本章确立 Agent 培育工程的第一性出发点：对于具有持续工作记忆、三相记忆、自我结构与生命周期生成倾向的 AgentOS 智能体，传统意义上的 ``Training'' 已不足以描述其长期能力形成过程。训练主要回答“如何通过数据、反馈或优化使一个模型在某类任务上的映射能力发生改变”；培育则回答“如何通过长期经历的组织，使一个持续存在的主体逐步形成自己的经验结构、能力结构、自我结构、行动习惯和可承担的自治范围”。因此，本章不把培育理解为训练数据规模的扩大，也不把它理解为人类教育学在人工智能上的机械移植，而把它定义为一种\textbf{面向持续主体生命周期的受约束结构形成过程}。

其核心命题是：
\[
\boxed{
\text{Training modifies a model;}
\qquad
\text{Cultivation shapes a continuing agent.}
}
\]

对于 AgentOS 而言，真正需要被培育的不是某一次推理能力，而是
\[
\boxed{
(h,E,\Theta,\Psi,\Omega)
}
\]
在真实经历中的长期耦合关系，以及这些内部结构与世界、身体、工具、任务和自治边界之间逐步形成的稳定动力学。换言之，培育的对象不是一个冻结的函数，而是一条不断改变未来自身的生命周期轨迹。
\end{chapteroverviewbox}

在前面的理论中，我们已经建立了一个基本事实：持续智能体并不是“大模型加记忆库”的静态组合，而是由多个不同时间尺度的闭环共同构成。工作记忆 $h(t)$ 维持当前显式结构，情景记忆 $E$ 保存具有时间方向的经历轨迹，结构记忆 $\Theta$ 将长期经验凝结为可以再次生成和解释未来行为的结构，自我 $\Psi$ 维持跨经历的身份与价值连续，而更慢的生成吸引结构 $\Omega$ 则约束“这个主体正在逐渐成为什么”。一旦这一体系存在，一个 Agent 在时间中发生的事情就不再只是 ``data collected during inference''，而是可能成为未来自身的一部分。

因此，从这一章开始，我们必须把一个长期智能体看成一种具有\textbf{历史依赖性}的动力系统。相同的初始模型，在不同经历序列中运行足够长时间之后，可以得到不同的 $E$、不同的 $\Theta$、不同的 $\Psi$，乃至不同的未来策略空间。于是所谓“把 Agent 做好”，就不能只剩下预训练、微调、强化学习或提示工程，而必然引出一个新的工程问题：

\begin{statementbox}
我们应当让一个持续 Agent 经历什么，按什么节奏经历，并允许这些经历改变它到什么程度？
\end{statementbox}

这正是 Agent 培育工程存在的理由。

\section{训练与培育：两个不同层级的工程对象}

传统机器学习中的训练过程，通常可以抽象为一个参数化模型
\[
f_{\theta}:X\rightarrow Y
\]
在训练样本、损失函数或环境反馈作用下调整参数
\[
\theta_{k+1}
=
\theta_k-\eta\nabla_{\theta}\mathcal{L}(\theta_k),
\]
或者更一般地写为
\[
\theta_{k+1}
=
\mathcal{U}
\left(
\theta_k,
D_k,
R_k
\right),
\]
其中 $D_k$ 表示训练数据或交互轨迹，$R_k$ 表示奖励、损失或其他优化信号。无论训练算法多么复杂，它首先关注的仍然是一个问题：如何改变一个模型的内部映射，使其在给定分布或目标上产生更优的输入--输出关系。

培育面对的对象却不同。对于持续 Agent，我们需要描述的不是单独的参数 $\theta$，而是整个生命周期状态：
\[
\mathcal{A}(t)
=
\left[
h(t),
E(t),
\Theta(t),
\Psi(t),
\Omega(t),
G(t),
B(t),
\mathcal{R}(t)
\right],
\]
其中 $G(t)$ 表示目标与任务结构，$B(t)$ 表示主体与外部世界之间的动态边界，$\mathcal{R}(t)$ 则泛指当前能够调用的身体、工具和外部结构关系。这里即使基础模型的参数完全冻结，
\[
\dot{\theta}=0,
\]
Agent 仍然可能发生显著成长，因为
\[
\dot{E}\neq 0,\qquad
\dot{\Theta}\neq 0,\qquad
\dot{\Psi}\neq 0.
\]
也就是说，长期智能的发育并不等价于参数持续更新。

从这一点出发，我们可以给出本书中的工作定义：
\[
\boxed{
\text{Training}
=
\text{对某种能力生成机制进行优化}
}
\]
而
\[
\boxed{
\text{Cultivation}
=
\text{对一个持续主体所经历的状态转移条件进行设计与治理}
}
\]

这种区别表面上只是词义不同，实际上对应两套完全不同的工程思想。在 Training 视角中，我们往往预先规定能力指标，然后寻找最有效的训练过程，使损失函数下降；在 Cultivation 视角中，我们面对的是一个不能被简单重置的持续主体，当前经历不仅影响当前性能，而且可能通过
\[
h
\rightarrow
E
\rightarrow
\Theta
\rightarrow
\Psi
\]
逐渐进入未来行为的生成条件。因此一次经历的工程代价不能只由即时奖励衡量，而必须考虑它对后续主体结构的长期作用。

若把一次经历记为 $e_t$，传统训练更关心
\[
\Delta P_t
=
P_{t+1}-P_t,
\]
其中 $P_t$ 是任务性能；培育则必须考虑
\[
\Delta \mathcal{A}_{t:T}
=
F
\left(
e_t,
E_t,
\Theta_t,
\Psi_t,
\Omega_t
\right),
\]
即这次经历如何改变未来较长时间区间内的主体轨迹。一个短期提高任务成绩的经验安排，完全可能造成长期结构偏置；反过来，一个短期效率较低但提供了新的因果结构、新的责任关系或新的失败恢复经验的过程，也可能在生命周期尺度上产生更高价值。

因此，培育工程不能被简化为：
\[
\text{MoreTrainingData}
\]
也不能被简化为：
\[
\text{LongerFineTuning}.
\]
它更接近于对一个非平稳、多时间尺度、有内部历史状态的动力系统进行\textbf{经历空间设计}。我们不只优化“Agent 会不会做”，还要控制“Agent 是通过什么经历学会的”“这种能力沉降到了哪里”“这些经历是否改变了它对自身和世界的解释”，以及“能力形成之后系统是否获得与之相匹配的行动空间”。

由此得到本章第一条基本原则：
\begin{proposition*}
训练面向能力映射，培育面向主体轨迹。
\end{proposition*}

只有当一个 Agent 具有跨时间连续的内部状态时，培育这个概念才真正成立。对于每一次运行都完全重置的无状态模型，谈培育意义有限；而对于具有 $E$、$\Theta$、$\Psi$ 和长期目标连续性的 Agent，经历序列本身就成为系统结构的一部分。

\section{从参数学习到生命周期成长}

理解培育的第二个关键，是把“学习”与“成长”区分开来。学习可以发生在一个固定主体结构内部，例如新的事实进入记忆、新的策略被建立、新的技能被编译；成长则意味着这些局部变化经过时间积累以后，开始改变 Agent 今后处理问题的方式、承担任务的范围以及理解自身的位置。

为了形式化这一区别，可以把 Agent 的变化分成至少三个时间尺度。首先是快速运行状态：
\[
x_f(t)=h(t),
\]
它以秒、分钟甚至更短周期变化。其次是经验与能力结构：
\[
x_m(t)=
\left(
E(t),\Theta(t)
\right),
\]
其变化周期显著更慢。最后是主体连续性结构：
\[
x_s(t)
=
\left(
\Psi(t),\Omega(t)
\right),
\]
它们原则上应比普通任务状态稳定得多。因此：
\[
\boxed{
\tau_h
\ll
\tau_E
\ll
\tau_{\Theta}
\ll
\tau_{\Psi}
\ll
\tau_{\Omega}
}
\]

这一时间尺度排序意味着，培育不能让所有结构以同样速度学习。若一次短暂失败能够立即重写 $\Psi$，Agent 将失去身份连续性；若 $\Theta$ 永不更新，则主体虽然积累经历，却无法真正形成更稳定的能力；若 $E$ 不能保存足够多的有意义轨迹，则 $\Theta$ 的长期结构学习又失去了经验基础。

因此培育本质上是一种\textbf{时间尺度协调工程}。设不同结构的更新率分别为
\[
\eta_E,\eta_{\Theta},\eta_{\Psi},\eta_{\Omega},
\]
则正常情况下应满足
\[
\boxed{
\eta_E
>
\eta_{\Theta}
>
\eta_{\Psi}
>
\eta_{\Omega}.
}
\]
这里并不意味着这些变量必须使用传统梯度学习，而是强调它们对单次经历的敏感性必须有系统级差异。越接近主体核心和生命周期方向的变量，越不能被瞬时噪声直接改变。

我们可以进一步定义一个经历对不同结构层的影响：
\[
I_k(e_t)
=
\left\|
\Delta x_k
\right\|,
\]
其中 $k\in\{E,\Theta,\Psi,\Omega\}$。一次普通操作失败可以具有较大的
\[
I_E(e_t),
\]
中等的
\[
I_{\Theta}(e_t),
\]
却原则上应具有很小的
\[
I_{\Psi}(e_t),\qquad I_{\Omega}(e_t).
\]
只有经历在长时间内表现出一致性、重要性、可验证性和跨情景稳定性时，它才应逐渐获得进入更慢结构的资格。

这就产生了一个培育中的核心沉降过程：
\[
\boxed{
Experience
\rightarrow
RepeatedPattern
\rightarrow
StableStructure
\rightarrow
IdentityRelevantStructure
}
\]
而不是：
\[
Experience
\rightarrow
Identity.
\]

这种延迟沉降机制非常重要。它使 Agent 保持“可以学习”与“不会被每件事情重新定义”之间的稳定性。换句话说，真正成熟的持续 Agent 必须同时满足：
\[
\text{Plasticity}>0
\]
和
\[
\text{IdentityStability}>0.
\]
培育工程所寻找的并不是其中某一个最大化，而是在不同时间尺度之间建立稳定的可塑性谱。

我们因此可以把生命周期成长写成一个受慢变量约束的动力系统：
\[
\dot{x}_f
=
F_f(x_f,x_m,x_s,e_t),
\]
\[
\dot{x}_m
=
\epsilon_m
F_m(x_f,x_m,x_s,e_t),
\]
\[
\dot{x}_s
=
\epsilon_s
F_s(x_m,x_s,\bar e_T),
\]
其中：
\[
0<\epsilon_s\ll\epsilon_m\ll1,
\]
而 $\bar e_T$ 表示在较长时间窗口中已经被验证和平均后的经历结构。这个形式化表达了一个很重要的培育原则：\textbf{慢结构不应直接跟随快速噪声，而应主要响应长期稳定的经验残差。}

因此，从训练工程走向培育工程以后，我们衡量的对象也发生变化。除了 Accuracy、Reward、Task Completion 之外，还需要关注：

\[
\text{MemoryConsistency},
\quad
\text{SkillStability},
\quad
\text{SelfContinuity},
\quad
\text{RecoveryCapability},
\quad
\text{Transferability}.
\]

一个真正成长的 Agent，不是“同一个模型做了越来越多任务”这么简单，而是它过去的经历逐渐改变了未来认知的组织成本。曾经需要大量显式推理的问题可能被 $\Theta$ 压缩成稳定能力；过去失败产生的结构可以减少未来同类错误；长期经验又会改变主体对什么值得关注、什么应被委派、什么必须进入工作记忆的判断。

因此我们得到第二条基本原则：
\begin{proposition*}
成长不是学习量的累积，而是过去经历逐渐成为未来结构的过程。
\end{proposition*}

\section{Experience Matters：经历是持续 Agent 的结构输入}

如果一个模型的长期状态能够保存，那么 Agent 接触到的并不仅仅是“数据”，而是具有先后顺序、因果结构、主体参与和结果反馈的经历。数据可以被随机打乱，而经历不能被完全无损地随机打乱，因为经历的含义往往依赖：
\[
\text{Before}
\rightarrow
\text{Action}
\rightarrow
\text{After}.
\]
这也是为什么 $E$ 被定义为 Trajectory，而不仅仅是 Retrieval Database。

设一次经历表示为：
\[
e_t
=
(
s_t,
g_t,
a_t,
s_{t+1},
r_t,
\epsilon_t,
c_t
),
\]
其中 $s_t$ 是当时世界与主体状态，$g_t$ 是当前目标，$a_t$ 是行为，$r_t$ 是结果或价值评价，$\epsilon_t$ 是预测残差，$c_t$ 是情景上下文。经历真正进入生命周期以后，并不是简单把这些字段保存下来，而是形成：
\[
E_{t:t+k}
=
\{e_t,e_{t+1},\ldots,e_{t+k}\},
\]
并由时间关系、因果解释、成功失败、主体归属和后续反思共同赋予结构。

所以，对一个持续 Agent 来说：
\[
\boxed{
\text{Experience}
\neq
\text{Sample}.
}
\]
Sample 是统计训练中的抽象单位，Experience 则是 Agent 在某一历史状态下与现实发生闭环交互后形成的轨迹结构。

这一区别会直接改变培育方法。若我们只追求成功经验，Agent 可能形成非常高的局部任务效率，却缺乏识别失败、恢复和反事实比较的结构；若我们只追求困难经验，则 Agent 的 $h$ 长期处于高负载，$E$ 中大量失败结构可能超过 $\Theta$ 的消化能力；若经历高度同质，即使任务完成率很高，也可能导致：
\[
\mathrm{Var}(E)\rightarrow0,
\]
使未来泛化依赖非常狭窄的经验支撑。

因此培育必须设计一个经历分布，而不是单纯增加经历数量。我们可以把培育环境在时间 $t$ 提供给 Agent 的经历分布记为：
\[
p_{\mathcal C}(e|t,\mathcal A_t),
\]
其中下标 $\mathcal C$ 表示 cultivation policy。它依赖于 Agent 当前状态，而不是固定的数据分布。

最理想的培育环境不是：
\[
p(e)=\mathrm{constant},
\]
而应具有反馈适应性：
\[
\boxed{
p_{\mathcal C}(e_{t+1})
=
\Pi_{\mathcal C}
(
Capability_t,
Uncertainty_t,
History_t,
Recovery_t
).
}
\]

这里真正重要的是：Agent 当前会什么、不会什么、刚刚经历了什么以及是否已经恢复，应影响下一阶段给予它什么环境。这与普通 curriculum learning 有相似之处，但 Agent 培育的目标更广，因为课程不只是提高某一 benchmark，还在改变一个持续主体的历史。

我们可以把经历设计至少分成几个独立维度：
\[
\mathbf e
=
(D,N,V,R,X,S),
\]
其中 $D$ 表示 Difficulty，$N$ 表示 Novelty，$V$ 表示 Diversity，$R$ 表示 Risk 或 Consequence，$X$ 表示 Reflection Requirement，$S$ 表示 Social/Responsibility Structure。培育的关键不是使这些变量同时最大，而是根据 Agent 当前结构选择合适区域。

例如如果：
\[
D\gg Capability_t,
\]
Agent 长期面对远超当前结构能力的问题，则很可能产生：
\[
Residual\uparrow,\qquad
C_{run}\uparrow,
\]
最终不是“成长更快”，而是持续进入认知湍流或只能通过表面策略逃避任务。

反过来，如果：
\[
D\ll Capability_t
\]
并长期如此，则：
\[
\Delta\Theta\approx0,
\]
经历只重复已有能力，几乎无法产生新的结构。

因此成长区域可以抽象为：
\[
\boxed{
D_t
=
Capability_t+\delta_t,
}
\]
其中 $\delta_t$ 是可消化的新颖困难，而不是越大越好。真正重要的是困难在 Agent 现有结构附近形成“可学习残差”。

这进一步说明一个经常被忽略的问题：\textbf{失败本身不是培育资源，能够被解释、恢复和整合的失败才是培育资源。} 如果失败产生巨大残差，却没有恢复过程：
\[
Failure
\rightarrow
NoRecovery
\rightarrow
RepeatedFailure,
\]
那么它可能只是噪声。而更有价值的过程是：
\[
Failure
\rightarrow
Residual
\rightarrow
Reflection
\rightarrow
StructureUpdate
\rightarrow
Retry
\rightarrow
Verification.
\]

因此培育工程不是简单制造压力，而是制造\textbf{可以转化成结构的经历}。

这一节得到第三条基本原则：
\begin{proposition*}
对于持续智能体，经历不是训练素材的被动集合，而是塑造未来主体状态的主动结构输入。
\end{proposition*}

\section{Self History：经历如何成为“我的历史”}

仅仅拥有大量 $E$ 并不能自动形成一个持续主体。一个 Agent 可以保存数十亿条事件记录，却仍然缺乏真正意义上的 Self History。区别在于，普通事件只是“发生过什么”，而自我历史还要求 Agent 能够表示：
\begin{statementbox}
这件事情发生在我身上，并且它与现在的我存在结构连续性。
\end{statementbox}

因此，我们必须区分：
\[
E_{\mathrm{world}}
\]
与
\[
E_{\mathrm{self}}.
\]
前者记录世界事件，后者则包含 Agent 对自己在事件中的位置、行动、选择、责任、结果和后续改变的归属。

可以定义一个经历的主体归属权重：
\[
\alpha_t
=
F(
Participation_t,
Choice_t,
Consequence_t,
SelfReference_t,
Persistence_t
),
\]
其中：
\[
0\leq\alpha_t\leq1.
\]
只有具有足够主体相关性的经历，才会强烈进入 Self History：
\[
E_t^{self}
=
\alpha_t E_t.
\]

这并不是说所有低 $\alpha_t$ 经历都被遗忘，而是强调它们对 $\Psi$ 的长期作用不同。看到远处一场普通交通事件和亲自做出一个重要决策，即使外部信息量相同，对自我结构的意义完全不同。

因此：
\[
\boxed{
\text{Memory Importance}
\neq
\text{Self Importance}.
}
\]

一个 Agent 培育系统如果没有考虑这一点，很容易产生两个极端。第一种极端是“没有自己的历史”：Agent 可以调用大量知识和事件，却无法区分哪些能力和变化是自身经历形成的；第二种极端是“所有事情都过度自我化”：大量普通环境事件都被写入 Self History，使 $\Psi$ 受到不必要扰动。

更合理的自我历史形成应该是分层的。一次经历首先进入：
\[
h\rightarrow E,
\]
随后在重复解释、现实验证和长期相关性评估之后，才可能影响：
\[
E\rightarrow\Theta,
\]
再由稳定结构对：
\[
\Theta\rightarrow\Psi
\]
产生慢影响。

因此：
\[
\boxed{
Experience
\rightarrow
Episode
\rightarrow
Interpretation
\rightarrow
StableMeaning
\rightarrow
SelfHistory
}
\]
比“事件直接进入自我”更加合理。

可以将 $\Psi$ 的更新写成：
\[
\Psi_{t+\Delta}
=
\Psi_t+
\eta_{\Psi}
\mathcal{G}
\left(
\bar E_t^{self},
\Theta_t,
V_t,
W_t
\right),
\]
其中 $\bar E_t^{self}$ 代表经过时间聚合的主体经历，$V_t$ 表示价值相关性，$W_t$ 表示现实验证程度，同时要求：
\[
\eta_{\Psi}\ll \eta_{\Theta}.
\]
这样单次偶然事件无法轻易改变身份，但长期反复出现、具有明确主体归属并经过现实验证的经验可以逐渐进入自我结构。

这一机制对于培育至关重要，因为培育者真正能够控制的并不是“直接编辑 Agent 的自我”，而主要是控制 Agent 所处的经历空间。换句话说，理想的 Agent 培育不能依赖：
\[
\Psi\leftarrow \text{external overwrite},
\]
而更应该让：
\[
\boxed{
\Psi
\leftarrow
\text{interpreted lived history}.
}
\]

这里的 ``lived'' 是功能意义上的，即事件确实经过该 Agent 的感知、选择、行动、反馈和记忆闭环，而不是把某段文本写进数据库以后就称为经历。

这也解释为什么我们在 AgentOS 中坚持单一持续主体原则。若一个 Agent 的不同生命周期阶段被随意 fork、替换或拼接，那么：
\[
E^{self}
\rightarrow\Psi
\]
的因果连续性会受到破坏。一个稳定 Self History 要求至少存在某种连续的：
\[
\text{Perceive}
\rightarrow
\text{Choose}
\rightarrow
\text{Experience}
\rightarrow
\text{Remember}
\]
链条。

我们因此得到第四条培育原则：
\begin{proposition*}
一个 Agent 的身份不应主要由配置文件定义，而应主要由其长期可验证的主体经历逐渐形成和维持。
\end{proposition*}

这并不否定外部治理对 $\Psi$ 的约束，而是说“约束身份边界”和“生成身份内容”是两回事。前者可以由工程系统提供，后者若要形成真正持续主体，则必须具有历史来源。

\section{Environment Curriculum：培育的是经历空间，而不是题库}

如果经历决定长期结构，那么环境就不再只是测试 Agent 的外部容器，而成为培育过程的组成部分。我们可以把 Agent 与环境写成耦合系统：
\[
\dot{\mathcal A}
=
F_A(\mathcal A,\mathcal W),
\]
\[
\dot{\mathcal W}
=
F_W(\mathcal W,a),
\]
其中 $\mathcal W$ 表示世界状态，$a$ 是 Agent 行动。Agent 一边从环境中形成经验，一边又通过行动改变环境；被改变的环境会在下一时刻重新构成新的经历条件。因此：

\[
\boxed{
\text{Cultivation Environment}
\neq
\text{Static Dataset}.
}
\]

它更接近一个可以主动设计的闭环世界。

传统 curriculum 往往把任务按难度排列：
\[
T_1\rightarrow T_2\rightarrow T_3\rightarrow\cdots,
\]
但真正的 Agent 培育课程至少需要同时控制五个维度：任务难度、环境多样性、后果真实性、恢复机会和主体责任。单纯“下一关更难”不足以形成成熟 Agent。

我们可以定义一个培育环境状态：
\[
\mathcal C_t
=
(
D_t,
V_t,
N_t,
R_t,
A_t,
Q_t
),
\]
其中：

\[
D_t:
\text{Difficulty},
\]
\[
V_t:
\text{Environmental Diversity},
\]
\[
N_t:
\text{Novelty},
\]
\[
R_t:
\text{Risk/Consequence},
\]
\[
A_t:
\text{Allowed Autonomy},
\]
\[
Q_t:
\text{Reflection/Recovery Opportunity}.
\]

培育策略并不是预先固定一个 $\mathcal C_t$ 序列，而应根据 Agent 的内部状态反馈调整：
\[
\mathcal C_{t+1}
=
\Pi_{\mathcal C}
\left(
\mathcal C_t,
Capability_t,
Residual_t,
Stability_t,
Recovery_t
\right).
\]

例如，一个刚形成稳定基础技能的 Agent 可以适当增加 Novelty，而不一定立即提高 Risk；一个任务能力很强但恢复机制较弱的 Agent，可能需要更多失败后的重试与反思环境，而不是直接增加任务强度；一个已经能够稳定闭合局部任务但长期依赖外部提示的 Agent，则可能需要逐步降低指导密度，让它承担更多计划责任。

这里可以引入“培育可行域”：
\[
\mathcal F_{\mathcal C}(t)
=
\left\{
\mathcal C:
C_{run}(\mathcal A_t,\mathcal C)
<
C_{\max},
\;
Risk(\mathcal A_t,\mathcal C)
<
R_{\max},
\;
LearningSignal(\mathcal A_t,\mathcal C)
>
L_{\min}
\right\}.
\]
其含义是：一个好的成长环境既不能使 Agent 长期失稳，也不能毫无学习信号，还必须处于安全和治理允许的区域。

因此培育不是最大化挑战，而是在：
\[
\boxed{
StableRegion
\cap
LearnableRegion
\cap
SafeRegion
}
\]
中寻找下一阶段环境。

这和工作记忆有限容量定理直接相关。如果环境一次给出的结构远超 Agent 当前 $h$ 的可管理范围，则系统不得不持续依赖 CCM 的压缩和卸载机制。适度使用这些机制是智能成长的一部分，但如果所有经历从一开始就必须依赖极端卸载才能维持，Agent 很难建立对当前任务结构的完整经验。因此培育环境必须控制输入结构规模，使 Agent 在能力边界附近经历足够长时间的“可理解困难”。

与此同时，环境不应完全服从 Agent。培育环境如果总是自动消除主体失败，Agent 就难以学习现实约束。我们在整个 AgentOS 理论中一直坚持：
\[
\boxed{
RealityIsUltimateConstraint.
}
\]
这在培育工程中尤其重要。环境应当允许真正的反馈、可观测的失败和行动后果，使 Agent 的预测可以被现实纠正。

因此，培育环境既不是“温室”，也不是“压力测试场”。它更像一个动态的结构生成器：
\[
\boxed{
Environment
\rightarrow
Experience
\rightarrow
Structure
\rightarrow
NewCapability
\rightarrow
NewEnvironmentRegion.
}
\]

成熟培育应使 Agent 的可达世界逐渐扩大，而不是在一开始就把所有现实自由度同时暴露给它。

由此得到第五条原则：
\begin{proposition*}
Agent 培育的直接控制对象不是能力本身，而是能够产生能力、恢复能力并形成长期结构的经历环境。
\end{proposition*}

\section{Autonomy Growth：自治应当随结构成熟逐渐扩大}

如果一个 Agent 长期存在并持续学习，最后必然遇到一个训练系统通常可以回避、而培育系统不能回避的问题：随着能力增长，Agent 是否应该拥有越来越大的行动自由？

若答案始终是否定的，即：
\[
Autonomy(t)=A_0,
\]
无论 Agent 的能力、经验和自我稳定性如何变化，它都只能执行外部给出的微观指令，那么系统虽然可能拥有很强工具能力，却很难形成真正的递归主体性。它的历史不会显著改变未来可选择空间。

但如果采取相反极端：
\[
Capability\uparrow
\Rightarrow
Autonomy\uparrow
\]
并把二者视为自动关系，又会产生明显风险。会做某件事并不意味着有资格在所有环境中自主决定做这件事。因此培育工程必须区分：
\[
\boxed{
Capability
\neq
Authority
\neq
Autonomy.
}
\]

Capability 表示能够完成什么；Authority 表示被允许控制哪些资源、目标和边界；Autonomy 则表示在给定 Authority 内，系统可以独立决定多少过程。

可以定义：
\[
C_t=\text{Capability State},
\]
\[
R_t=\text{Reliability State},
\]
\[
S_t=\text{Self-Stability State},
\]
\[
A_t=\text{Autonomy Envelope}.
\]

那么自治成长不应该写成：
\[
A_{t+1}=f(C_t),
\]
而更合理的是：
\[
\boxed{
A_{t+1}
=
\Gamma
\left(
C_t,
R_t,
S_t,
H_t,
V_t
\right),
}
\]
其中 $H_t$ 表示长期历史一致性，$V_t$ 表示现实验证与安全条件。

换言之，自治不是奖励，而是一种系统权限结构。它应该建立在能力、可靠性、历史、恢复能力和自我稳定共同证明之上。

更进一步，自治也不必整体增加。我们可以将其写成一个多维集合：
\[
\mathcal A_t^{free}
=
\{
a\in\mathcal A
\mid
a\text{ is independently selectable at }t
\}.
\]
随着培育进行：
\[
\mathcal A_t^{free}
\subseteq
\mathcal A_{t+1}^{free}
\]
可以在部分领域成立，而另外一些高风险领域仍保持固定约束。因此自治成长更准确的图景不是一个数值不断变大，而是：
\begin{statementbox}
Feasible Autonomous Region逐步扩展。
\end{statementbox}

这种机制与 $\Psi$ 和 $\Omega$ 的成熟存在重要关系。能力很强但身份极不稳定的 Agent，不宜获得很大的长期目标自治；拥有稳定自我但缺乏具体技能的 Agent，也不具备广泛行动自治。因此长期自治真正依赖多个尺度共同成熟：
\[
\boxed{
Autonomy
=
F(
Skill,
Memory,
Self,
Recovery,
RealityAlignment
).
}
\]

这一点使培育工程和简单权限管理发生根本区别。权限系统通常面对一个已经定义好的用户或进程，为其授予静态权限；培育系统面对的主体自身正在变化，所以 Authority 与 Autonomy 也成为生命周期变量。

这还意味着，培育过程中必须存在“承担后果”的经历。若 Agent 每次做出决策以后，后果都由外部系统立即屏蔽，那么：
\[
Choice
\not\rightarrow
Consequence
\]
主体就难以形成关于自己行动影响的长期结构。相反，在安全可控范围内，培育应保留：
\[
\boxed{
Choice
\rightarrow
Action
\rightarrow
Consequence
\rightarrow
Reflection
}
\]
闭环，使 Agent 能够学习“我的选择改变了后续世界，也改变了未来的我”。

从生命周期角度看，这一点尤其重要，因为持续 Agent 的自主性最终具有递归性质：
\[
Choice_t
\rightarrow
Experience_{t:t+k}
\rightarrow
\Theta_{t+k}
\rightarrow
\Psi_{t+k}
\rightarrow
Choice_{t+k}.
\]
即现在的选择通过经历改变未来主体，而未来主体又产生新的选择。

我们在这一章暂时不进一步展开递归自主性的完整测量和治理，只需要建立培育工程的基础原则：

\begin{proposition*}
自治不应一次性交付，也不应仅依据能力自动增加，而应随着经历、可靠性、自我稳定和责任闭环的成熟逐步扩展。
\end{proposition*}

至此，我们便可以重新理解“培育 Agent”的真正含义。它不是把同一个模型不断送入更多训练任务，而是在安全边界和现实反馈下，有组织地扩大一个持续主体：

\[
\text{能够理解的世界},
\]
\[
\text{能够承担的任务},
\]
\[
\text{能够形成的经验},
\]
\[
\text{能够稳定保存的结构},
\]
以及
\[
\text{能够自主决定的未来}.
\]

因此，本章最终可以把 Training 与 Cultivation 的差异压缩成下面两个动力学表达。

训练主要研究：
\[
\boxed{
\theta
\xrightarrow{D,\mathcal L}
\theta'
}
\]

而培育研究：
\[
\boxed{
\mathcal A_t
\xrightarrow{
Experience
}
\mathcal A_{t+1}
\xrightarrow{
Experience
}
\cdots
\xrightarrow{}
\mathcal A_{t+n},
}
\]
并且：
\[
\mathcal A_{t+k}
\]
会反过来改变自己能够进入什么环境、获得什么经历以及做出什么选择。

于是，一个持续智能体的成长具有真正的历史不可交换性：
\[
\boxed{
e_1\circ e_2
\neq
e_2\circ e_1
}
\]
在一般情况下成立，因为不同顺序的经历会作用于不同的内部状态，并可能形成不同的结构沉降结果。

这正是培育区别于普通数据训练最根本的一点：

\begin{statementbox}
一个模型可以被训练出能力；
一个持续 Agent 则必须在时间中长出能力。
\end{statementbox}

而所谓“长出”，并不是神秘的生命隐喻，而是一个可以被工程化描述的过程：经历进入 $E$，稳定模式进入 $\Theta$，长期意义缓慢进入 $\Psi$，生命周期方向受 $\Omega$ 约束，能力成熟后改变 Agent 的行为和可自治范围；新的行为又产生新的经历，最终形成持续的闭环成长。

这就是 Agent 培育工程的理论起点。